\documentclass[12pt,a4paper]{report}
\usepackage[utf8]{inputenc}
\usepackage{amsmath}
\usepackage{amsfonts}
\usepackage{amssymb}
\usepackage{amsthm}
\usepackage{hyperref}

\usepackage{multicol}
\usepackage{fancyhdr}
\usepackage[inline]{enumitem}
\usepackage{tikz}
\usepackage{tikz-cd}
\usetikzlibrary{calc}
\usetikzlibrary{shapes.geometric}
\usepackage[margin=0.5in]{geometry}
\usepackage{xcolor}
\usepackage{ esint }

\hypersetup{
    colorlinks=true,
    linkcolor=blue,
    filecolor=magenta,      
    urlcolor=cyan,
    pdftitle={Functional Analysis},
    pdfpagemode=FullScreen,
    }

%\urlstyle{same}

\newcommand{\CLASSNAME}{Comprehensive Exams -- Functional Analysis}
\newcommand{\STUDENTNAME}{Paul Carmody}
\newcommand{\ASSIGNMENT}{NOTES }
\newcommand{\DUEDATE}{April 2026}
\newcommand{\SEMESTER}{Spring 2026}
\newcommand{\SCHEDULE}{TBD}
\newcommand{\ROOM}{Crawford 332}

\newcommand{\MMN}{M_{m\times n}}
\newcommand{\FF}{\mathcal{F}}

\pagestyle{fancy}
\fancyhf{}
\chead{ \fancyplain{}{\CLASSNAME} }
%\chead{ \fancyplain{}{\STUDENTNAME} }
\rhead{\thepage}
\fancyfoot[C]{-- \thepage --}
\input{../../MyMacros}

\newcommand{\RED}[1]{\textcolor{red}{#1}}
\newcommand{\BLUE}[1]{\textcolor{blue}{#1}}
\newcommand{\FFF}{\mathcal{F}}

\begin{document}

\begin{center}
	\Large{\CLASSNAME -- \SEMESTER} 
\end{center}
\begin{center}
	\STUDENTNAME \\
	\ASSIGNMENT -- \DUEDATE\\
\end{center} 

\section*{Functional Analysis -- Overview}

Functional analysis is a broad subject, but it has a very coherent ``architecture." Here’s a structured outline that shows how the pieces fit together and what each part is trying to accomplish.

---
\begin{description}

\item \textbf{1. Foundations: Normed and Metric Spaces}

\item \textbf{Goal:} Generalize notions of distance, convergence, and continuity.

\begin{itemize}
    \item \DEFINE{Metric spaces}

	\begin{itemize}
        \item Open/closed sets, completeness
        \item Compactness, connectedness
	\end{itemize}

    \item \DEFINE{Normed vector spaces}

	\begin{itemize}
        \item Norms, induced metrics
        \item Examples: $ \mathbb{R}^n $, $ C([a,b]$ ), $ L^p $
	\end{itemize}

    \item \DEFINE{Banach spaces (complete normed spaces)}

	\begin{itemize}
        \item Importance of completeness
        \item Examples: $ L^p $, $ C([a,b]$ )
	\end{itemize}

\end{itemize}
---

\item \textbf{2. Hilbert Spaces}

\item \textbf{Goal:} Introduce geometry (angles, orthogonality) into infinite dimensions.

\begin{itemize}
    \item Inner product spaces
    \item Hilbert spaces (complete inner product spaces)
    \item Orthogonality and projections
    \item Orthonormal bases
    \item Fourier series and expansions
\end{itemize}

---

\item \textbf{3. Linear Operators}

\item \textbf{Goal:} Study functions between infinite-dimensional spaces.

\begin{itemize}
    \item Bounded (continuous) linear operators

    \item Operator norm

    \item Examples:

	\begin{itemize}
        \item Integral operators
        \item Differential operators
	\end{itemize}

    \item Finite vs infinite dimensional behavior
\end{itemize}

---

\item \textbf{4. Duality and Functionals}

\item \textbf{Goal:} Understand spaces via linear functionals.

\begin{itemize}
    \item Dual space $ X^* $
    \item Hahn–Banach theorem (extension of functionals)
    \item Weak and weak-* convergence
\end{itemize}

---

\item \textbf{5. Core Theorems of Functional Analysis}

These are the ``big three" structural results:

\begin{itemize}
    \item \DEFINE{Hahn–Banach Theorem}

	\begin{itemize}
        \item Extension of linear functionals
	\end{itemize}
    \item \DEFINE{Uniform Boundedness Principle (Banach–Steinhaus)}

	\begin{itemize}
        \item Control of operator families
	\end{itemize}
    \item \DEFINE{Open Mapping Theorem}

	\begin{itemize}
        \item Surjective bounded operators are open
	\end{itemize}
    \item \DEFINE{Closed Graph Theorem}
\end{itemize}

These results rely heavily on completeness.

---

\item \textbf{6. Topologies in Infinite Dimensions}

\item \textbf{Goal:} Introduce weaker notions of convergence.

\begin{itemize}
    \item Weak topology
    \item Weak-* topology
    \item Compactness (e.g., Banach–Alaoglu theorem)
\end{itemize}

---

\item \textbf{7. Spectral Theory}

\item \textbf{Goal:} Generalize eigenvalues and diagonalization.

\begin{itemize}
    \item Spectrum of an operator
    \item Resolvent
    \item Spectral radius
    \item Compact operators
    \item Self-adjoint operators (Hilbert spaces)
    \item Spectral theorem
\end{itemize}

---

\item \textbf{8. Operator Algebras}

\item \textbf{Goal:} Study collections of operators as algebraic objects.

\begin{itemize}
    \item Banach algebras
    \item $ C^* $-algebras
    \item Applications to quantum mechanics
\end{itemize}

---

\item \textbf{9. Distributions and Sobolev Spaces}

\item \textbf{Goal:} Extend notion of functions to handle PDEs.

\begin{itemize}
    \item Test functions and distributions
    \item Weak derivatives
    \item Sobolev spaces $ W^{k,p} $
    \item Embedding theorems

\end{itemize}

---

\item \textbf{10. Applications}

Functional analysis is largely motivated by applications:

\begin{itemize}
    \item \DEFINE{Partial Differential Equations}

	\begin{itemize}
        \item Weak solutions
        \item Variational methods
	\end{itemize}
    \item \DEFINE{Quantum mechanics}

	\begin{itemize}
        \item Hilbert space framework
	\end{itemize}
    \item \DEFINE{Harmonic analysis}

	\begin{itemize}
        \item Fourier transform
	\end{itemize}
    \item \DEFINE{Numerical analysis}
    \item \DEFINE{Optimization and control}
\end{itemize}

---

\item \textbf{11. Advanced Topics (Optional Extensions)}

\begin{itemize}
    \item Semigroup theory (evolution equations)
    \item Nonlinear functional analysis
    \item Fixed point theorems (Banach, Schauder)
    \item Fredholm theory
    \item Measure theory connections
\end{itemize}

---

\item \textbf{Big Picture}

Functional analysis is essentially:

\DEFINE{Linear algebra + topology + limits in infinite dimensions}

It studies:

\begin{itemize}
    \item Spaces of functions
    \item Operators acting on them
    \item Convergence and structure
\end{itemize}

---
\end{description}

If you want, I can tailor this outline specifically toward \DEFINE{PDEs}, \DEFINE{geometry}, or a \DEFINE{first graduate course}, which each emphasize very different parts.

\newpage

\subsection*{1. Foundations: Normed and Metric Spaces}

Here’s a focused ``toolkit-style" outline of the \DEFINE{core ideas, definitions, and theorems} for each of the three layers:

---

\begin{description}
\item \textbf{1. Metric Spaces}

    \begin{itemize}
    \item \textbf{Key Idea}

Abstract the notion of \DEFINE{distance} to study convergence and topology without linear structure.

    \item \textbf{Core Definitions}

            \begin{itemize}
            \item \textit{\DEFINE{Metric space}: $(X,d)$ with}
\begin{align*}
  d(x,y)\ge 0,\quad d(x,y)=0 \iff x=y,\quad d(x,y)=d(y,x),\quad d(x,z)\le d(x,y)+d(y,z)
\end{align*}

            \item \textit{\DEFINE{Open ball}: $B_r(x)={y\in X: d(x,y)<r}$}

            \item \textit{\DEFINE{Open / closed sets}}

            \item \textit{\DEFINE{Closure}, \DEFINE{interior}, \DEFINE{boundary}}

            \item \textit{\DEFINE{Convergence}: $x_n \to x \iff d(x_n,x)\to 0$}

            \item \textit{\DEFINE{Cauchy sequence}:}
\begin{align*}
  d(x_n,x_m)\to 0 \quad \text{as } n,m\to\infty
\end{align*}

            \item \textit{\DEFINE{Complete space}: every Cauchy sequence converges}
            \end{itemize}

---

    \item \textbf{Key Concepts}

            \begin{itemize}
            \item \textit{Continuity via metric}
            \item \textit{Sequential vs topological definitions}
            \item \textit{Compactness vs completeness}
            \end{itemize}

---

    \item \textbf{Fundamental Theorems}

            \begin{itemize}
            \item \textit{\DEFINE{Triangle inequality consequences}}

  * Reverse triangle inequality:
    \begin{align*}
    |d(x,z)-d(y,z)| \le d(x,y)
    \end{align*}

            \item \textit{\DEFINE{Characterization of continuity}}

  * Equivalent to sequential continuity

            \item \textit{\DEFINE{Compactness (metric spaces)}}
  Equivalent:

  * Sequential compactness
  * Every sequence has convergent subsequence
  * Open cover definition

            \item \textit{\DEFINE{Heine–Borel (in $\mathbb{R}^n$)}}

  * Compact $\iff$ closed and bounded

            \item \textit{\DEFINE{Cantor Intersection Theorem}}

  * Nested closed sets with diameters $\to$ 0 have single-point intersection

            \end{itemize}
---

    \item \textbf{Important Lemmas}

            \begin{itemize}
            \item \textit{Limit in metric space is unique}
            \item \textit{Cauchy $\implies$ bounded}
            \item \textit{Convergent $\implies$ Cauchy}
            \end{itemize}
    \end{itemize}

---

\item \textbf{2. Normed Vector Spaces}

    \begin{itemize}
    \item \textbf{Key Idea}

Add \DEFINE{linear structure + size (norm)} $\to$ geometry + analysis.

---

    \item \textbf{Core Definitions}

            \begin{itemize}
            \item \textit{\DEFINE{Normed space} $(X,|\cdot|)$:}
\begin{align*}
  |x|\ge 0,\quad |x|=0 \iff x=0,\quad |\lambda x|=|\lambda||x|,\quad |x+y|\le |x|+|y|
\end{align*}

            \item \textit{Induced metric:}
\begin{align*}
  d(x,y)=|x-y|
\end{align*}

            \item \textit{\DEFINE{Convergence}: $|x_n-x|\to 0$}

            \item \textit{\DEFINE{Linear operator} $T:X\to Y$}

            \item \textit{\DEFINE{Bounded operator}:}
\begin{align*}
  |Tx|\le C|x|
\end{align*}

            \item \textit{\DEFINE{Operator norm}:}
\begin{align*}
  |T|=\sup_{|x|\le 1}|Tx|
\end{align*}

            \end{itemize}
---

    \item \textbf{Key Concepts}

            \begin{itemize}
            \item \textit{Geometry (balls, convexity)}
            \item \textit{Linear structure + topology interaction}
            \item \textit{Finite vs infinite dimensional differences}
            \end{itemize}

---

    \item \textbf{Fundamental Theorems}

            \begin{itemize}
            \item \textit{\DEFINE{Equivalence of norms (finite dimension)}}

  * All norms on finite-dimensional spaces are equivalent

            \item \textit{\DEFINE{Continuity of linear maps}}
  Equivalent:

            \begin{itemize}
                 \item   Continuous
                 \item   Continuous at 0
                 \item   Bounded
            \end{itemize}

            \item \textit{\DEFINE{Riesz lemma}}

            \begin{itemize}
                 \item In infinite-dimensional spaces, unit ball is ``large" 
                 \item Used to show non-compactness
            \end{itemize}

            \item \textit{\DEFINE{Hahn–Banach Theorem (normed version)}}

  * Extend bounded linear functionals without increasing norm

            \end{itemize}
---

    \item \textbf{Important Lemmas / Results}

            \begin{itemize}
            \item \textit{\DEFINE{Triangle inequality variants}}
            \item \textit{\DEFINE{Convexity of norm balls}}
            \item \textit{\DEFINE{Bounded sets vs totally bounded sets}}
            \end{itemize}

    \end{itemize}
---

\item \textbf{3. Banach Spaces}

    \begin{itemize}
    \item \textbf{Key Idea}

A \DEFINE{complete normed space} $\to$ limits behave well $\to$ powerful global theorems.

---

    \item \textbf{Core Definition}

            \begin{itemize}
            \item \textit{\DEFINE{Banach space}: normed space where every Cauchy sequence converges}

            \end{itemize}
---

    \item \textbf{Key Concepts}

            \begin{itemize}
            \item \textit{Completeness enables:}

  * Existence theorems
  * Stability of limits
  * Functional analytic structure

            \end{itemize}
---

    \item \textbf{Fundamental Theorems (the ``Big Four")}

		\begin{description}
        \item \textbf{1. Hahn–Banach Theorem}

            \begin{itemize}
            \item \textit{Extend linear functionals:}
\begin{align*}
  f: Y \to \mathbb{R} \quad \Rightarrow \quad \tilde f: X \to \mathbb{R}
\end{align*}
            \item \textit{Preserves norm}

\DEFINE{Impact:}

            \item \textit{Rich dual spaces}
            \item \textit{Separation of points and sets}
            \end{itemize}

---

        \item \textbf{2. Uniform Boundedness Principle (Banach–Steinhaus)}

If $ {T_n} \subset \mathcal{L}(X,Y) $ and
\begin{align*}
\sup_n |T_n x| < \infty \quad \forall x \in X
\end{align*}
then
\begin{align*}
\sup_n |T_n| < \infty
\end{align*}

\DEFINE{Idea:} pointwise bounded $\implies$ uniformly bounded

---

        \item \textbf{3. Open Mapping Theorem}

If $T:X\to Y$ is bounded, linear, surjective (Banach spaces), then:

            \begin{itemize}
            \item \textit{$T$ maps open sets to open sets}
            \end{itemize}

---

        \item \textbf{4. Closed Graph Theorem}

If $T:X\to Y$ has closed graph, then $T$ is bounded

		\end{description}
---

    \item \textbf{Other Key Results}

            \begin{itemize}
            \item \textit{\DEFINE{Inverse Mapping Theorem}}

  * Bijective bounded linear map $\implies$ inverse is bounded

            \item \textit{\DEFINE{Banach Fixed Point Theorem (Contraction Mapping)}}

  * Contractions have unique fixed points

            \item \textit{\DEFINE{Baire Category Theorem}}

  * Complete metric spaces are ``large"
  * Foundation of many major results

            \end{itemize}
---

    \item \textbf{Important Concepts}

            \begin{itemize}
            \item \textit{Dense subsets}
            \item \textit{Nowhere dense sets}
            \item \textit{Category (meager vs residual)}
            \end{itemize}

    \end{itemize}
---

\item \textbf{Big Structural Picture}

You can think of the progression as:

        \begin{itemize}
        \item \textbf{Metric spaces} $\to$ convergence + topology

        \item \textbf{Normed spaces} $\to$ add linear structure

        \item \textbf{Banach spaces} $\to$ add completeness $\to$ \DEFINE{powerful global theorems}

        \end{itemize}
---

\item \textbf{What to Remember Most}

If you’re prioritizing:

    \begin{itemize}
        \item \textbf{Definitions to know cold}

            \begin{itemize}
            \item \textit{Metric, norm, completeness, Cauchy}
            \item \textit{Bounded linear operator}
            \end{itemize}

        \item \textbf{Theorems that drive everything}

            \begin{itemize}
            \item \textit{Hahn–Banach}
            \item \textit{Uniform Boundedness}
            \item \textit{Open Mapping}
            \item \textit{Closed Graph}
            \item \textit{Banach Fixed Point}
            \item \textit{Baire Category}
            \end{itemize}
    \end{itemize}

---
\end{description}

If you want, I can next:

            \begin{itemize}
            \item Give \DEFINE{typical exam problems} from each section, or
            \item Show how these results are used in \DEFINE{PDEs or Sobolev spaces} (which connects to your earlier questions).
            \end{itemize}

\newpage
\subsection*{Hilbert Spaces}

Here is a \DEFINE{structured ``core toolkit" for Hilbert spaces}—focusing on the ideas, definitions, and theorems that actually drive the subject.

---

\begin{description}
\item \DEFINE{1. Core Idea of Hilbert Spaces}

A Hilbert space is:

$\to$ A \DEFINE{complete inner product space} → combines geometry (angles, orthogonality) with analysis (limits).

It is the infinite-dimensional analogue of \DEFINE{Euclidean space}.

---

\item \DEFINE{2. Fundamental Definitions}

\begin{description}
    \item \textbf{Inner Product}

A function $ \langle \cdot,\cdot \rangle : H \times H \to \mathbb{C} $ (or $\mathbb{R}$) such that:

\begin{itemize}
\item Linearity (in first argument)
\item Conjugate symmetry:
  \begin{align*}
  \langle x,y\rangle = \overline{\langle y,x\rangle}
  \end{align*}
\item Positive definiteness:
  \begin{align*}
  \langle x,x\rangle \ge 0,\quad =0 \iff x=0
  \end{align*}

---
\end{itemize}

    \item \textbf{Norm induced by inner product}

\begin{align*}
|x| = \sqrt{\langle x,x\rangle}
\end{align*}

---

    \item \textbf{Hilbert Space}

A space $H$ that is:

\begin{itemize}
\item Inner product space
\item Complete with respect to $|\cdot|$
\end{itemize}

---

    \item \textbf{Orthogonality}

\begin{align*}
x \perp y \iff \langle x,y\rangle = 0
\end{align*}

---

    \item \textbf{Orthonormal Set}

\begin{align*}
\langle e_i, e_j\rangle = \delta_{ij}
\end{align*}

---
\end{description}

\item \DEFINE{3. Key Geometric Ideas}

\begin{itemize}
\item Angles and projections make sense
\item Best approximation = orthogonal projection
\item Infinite-dimensional ``Euclidean geometry"
\end{itemize}

---

\item \DEFINE{4. Fundamental Inequalities}

\begin{description}
    \item \textbf{Cauchy–Schwarz Inequality}

\begin{align*}
|\langle x,y\rangle| \le |x||y|
\end{align*}

    \item \textbf{Triangle Inequality}

Follows from Cauchy–Schwarz

    \item \textbf{Parallelogram Law}

\begin{align*}
|x+y|^2 + |x-y|^2 = 2|x|^2 + 2|y|^2
\end{align*}

\DEFINE{Important:} Characterizes inner product norms.

---
\end{description}

\item \DEFINE{5. Orthogonality and Projections}

\begin{description}
    \item \textbf{Orthogonal Complement}

For $M \subset H$:
\begin{align*}
M^\perp = {x \in H : \langle x,m\rangle = 0 \ \forall m\in M}
\end{align*}

---

    \item \textbf{Projection Theorem (CRUCIAL)}

If $M$ is a \DEFINE{closed subspace} of $H$, then:

For every $x\in H$, there exists a unique $y\in M$ such that
\begin{align*}
|x-y| = \inf_{m\in M} |x-m|
\end{align*}

And:
\begin{align*}
x - y \perp M
\end{align*}

---

    \item \textbf{Decomposition}

\begin{align*}
H = M \oplus M^\perp
\end{align*}

---
\end{description}

\item \DEFINE{6. Orthonormal Bases and Expansions}

\begin{description}
    \item \textbf{Orthonormal Basis (ONB)}

A set ${e_i}$ such that:

* Orthonormal
* Span is dense in $H$

---

    \item \textbf{Fourier Expansion}

For ONB ${e_i}$:
\begin{align*}
x = \sum \langle x,e_i\rangle e_i
\end{align*}

---

    \item \textbf{Parseval’s Identity}

\begin{align*}
|x|^2 = \sum |\langle x,e_i\rangle|^2
\end{align*}

---

    \item \textbf{Bessel’s Inequality}

\begin{align*}
\sum |\langle x,e_i\rangle|^2 \le |x|^2
\end{align*}

---
\end{description}

\item \DEFINE{7. Key Structural Theorems}

\begin{description}
    \item \textbf{Riesz Representation Theorem (VERY IMPORTANT)}

Every bounded linear functional $f \in H^*$ has the form:
\begin{align*}
f(x) = \langle x,y\rangle
\quad \text{for a unique } y \in H
\end{align*}

\DEFINE{Meaning:}
\begin{align*}
H^* \cong H
\end{align*}

---

    \item \textbf{Gram–Schmidt Process}

* Converts linearly independent set → orthonormal set

---

    \item \textbf{Pythagorean Theorem}

If $x \perp y$:
\begin{align*}
|x+y|^2 = |x|^2 + |y|^2
\end{align*}

---
\end{description}

\item \DEFINE{8. Operators on Hilbert Spaces}

\begin{description}
    \item \textbf{Bounded Linear Operators}

Same as Banach space theory, but with extra structure.

---

    \item \textbf{Adjoint Operator}

For $T:H\to H$, define $T^*$ by:
\begin{align*}
\langle Tx,y\rangle = \langle x,T^*y\rangle
\end{align*}

---

    \item \textbf{Special Classes}

\begin{itemize}
\item \DEFINE{Self-adjoint}: $T = T^*$
\item \DEFINE{Unitary}: $T^*T = I$
\item \DEFINE{Normal}: $TT^* = T^*T$
\end{itemize}

---
\end{description}

\item \DEFINE{9. Spectral Theory (Hilbert Version)}

\begin{description}
    \item \textbf{Spectrum}

Generalization of eigenvalues

---

    \item \textbf{Spectral Theorem (KEY RESULT)}

For self-adjoint (or normal) operators:

* Operator can be ``diagonalized" in an appropriate sense

Finite-dimensional analogy:

* Symmetric matrix → orthonormal eigenbasis

Infinite-dimensional:

* Integral decomposition / spectral measures

---
\end{description}

\item \DEFINE{10. Weak Convergence}

\begin{description}
    \item \textbf{Weak Convergence}

\begin{align*}
x_n \rightharpoonup x \iff \langle x_n, y\rangle \to \langle x,y\rangle \ \forall y\in H
\end{align*}

---

    \item \textbf{Key Facts}

\begin{itemize}
\item Weak convergence + norm convergence differ
\item Bounded sequences have weakly convergent subsequences (in many cases)
\end{itemize}

---
\end{description}

\item \DEFINE{11. Important Lemmas and Facts}

\begin{itemize}
\item Orthonormal sets are linearly independent
\item Closed subspaces are Hilbert spaces
\item $ (M^\perp)^\perp = \overline{M} $
\item Finite-dimensional subspaces are closed
\end{itemize}

---

\item \DEFINE{12. Canonical Examples}

\begin{itemize}
\item $ \ell^2 $: square-summable sequences
\item $ L^2(\Omega$ ): square-integrable functions
\item $ \mathbb{R}^n, \mathbb{C}^n $
\end{itemize}

---

\item \DEFINE{Big Picture}

Hilbert spaces give you:

\begin{description}
\item \textbf{Geometry}

* Angles, orthogonality, projections

\item \textbf{Analysis}

* Convergence, completeness

\item \textbf{Algebra of operators}

* Adjoint, spectral theory

\end{description}
---

\item \DEFINE{What to Know Cold}

If you’re prioritizing:

\begin{description}
\item \textbf{Definitions}

* Inner product, orthogonality, ONB

\item \textbf{Core Theorems}

* Cauchy–Schwarz \\
* Projection theorem \\
* Riesz representation \\
* Parseval identity \\
* Spectral theorem (at least conceptually) \\

\end{description}
---

\end{description}

If you want, I can next:

* Compare \DEFINE{Hilbert vs Banach spaces (what fails without inner products)}, or \\
* Work through \DEFINE{typical proofs (e.g., Riesz, projection theorem)} step-by-step.



\end{document}
